% Individual Reflection — Domi (individual submission).
% Build: latexmk -pdf individual-reflection-domi.tex
\documentclass[11pt,a4paper]{article}
% Shared style for the StudyOS Thesis-Finder course submissions.
% Included by every submission document via % Shared style for the StudyOS Thesis-Finder course submissions.
% Included by every submission document via % Shared style for the StudyOS Thesis-Finder course submissions.
% Included by every submission document via \input{preamble}.
% Plain black-and-white house style: no colour, prints identically in mono.

\usepackage[T1]{fontenc}
\usepackage[utf8]{inputenc}
\usepackage{lmodern}
\usepackage[english]{babel}
\usepackage{microtype}
\usepackage[a4paper,margin=2.4cm]{geometry}
\usepackage{xcolor}
\usepackage{titlesec}
\usepackage{enumitem}
\usepackage{booktabs}
\usepackage{tabularx}
\usepackage{tcolorbox}
\tcbuselibrary{breakable,skins}
\usepackage{pifont}
\usepackage{xurl}
\usepackage[hidelinks]{hyperref}

\urlstyle{tt}

% --- Palette (plain black and white) ----------------------------------------
% Everything is black on white; the only greys are for a hairline rule and the
% callout title bar, so the documents print identically in mono.
\definecolor{sosAccent}{gray}{0}
\definecolor{sosDark}{gray}{0}
\definecolor{sosOpen}{gray}{0}
\definecolor{sosOpenBg}{gray}{1}
\definecolor{sosRule}{gray}{0.6}

% --- Headings ---------------------------------------------------------------
\titleformat{\section}{\normalfont\large\bfseries}{\thesection}{0.6em}{}
\titleformat{\subsection}{\normalfont\normalsize\bfseries\itshape}{\thesubsection}{0.6em}{}
\titlespacing*{\section}{0pt}{1.1\baselineskip}{0.35\baselineskip}
\titlespacing*{\subsection}{0pt}{0.7\baselineskip}{0.2\baselineskip}

\setlist[itemize]{leftmargin=1.2em,itemsep=0.15em,topsep=0.25em,parsep=0pt}
\setlist[enumerate]{leftmargin=1.4em,itemsep=0.15em,topsep=0.25em,parsep=0pt}

% --- Semantic helpers -------------------------------------------------------
% Inline reference to a file or path inside the project repository.
\newcommand{\repofile}[1]{\path{#1}}

% One "the assignment asks for ..." lead-in line.
\newcommand{\asks}[1]{\noindent\textbf{Assignment asks for:} #1\par\vspace{0.3em}}

% Open question the team still has to answer (was <mark>highlighted</mark> in the
% markdown draft).
\newenvironment{openbox}[1][Open --- needs the team]{%
  \begin{tcolorbox}[
    breakable,
    colback=white,
    colframe=black,
    colbacktitle=black!8,
    coltitle=black,
    boxrule=0.4pt,
    titlerule=0pt,
    left=0.6em,right=0.6em,top=0.5em,bottom=0.5em,
    title={#1},
    fonttitle=\bfseries\small,
    sharp corners,
  ]\small
}{%
  \end{tcolorbox}
}

\NewDocumentEnvironment{infobox}
  { O{Open --- needs the team} O{} }
{%
  \begin{tcolorbox}[
    breakable,
    colback=white,
    colframe=black,
    colbacktitle=black!8,
    coltitle=black,
    boxrule=0.4pt,
    titlerule=0pt,
    left=0.6em,
    right=0.6em,
    top=0.5em,
    bottom=0.5em,
    title={#1},
    fonttitle=\bfseries\small,
    sharp corners,
    #2
  ]\small
}{%
  \end{tcolorbox}
}


% Status markers for the submission checklist.
\newcommand{\stdone}{\ding{51}}
\newcommand{\stpartial}{\ding{108}}
\newcommand{\stgap}{\ding{55}}

% Title block.
\newcommand{\sostitle}[3]{%
  \begin{center}
    {\LARGE\bfseries #1}\\[0.35em]
    {\large\itshape #2}\\[0.5em]
    {\small #3}
  \end{center}
  \vspace{0.4em}
  {\color{sosRule}\hrule height 0.4pt}
  \vspace{0.9em}
}
.
% Plain black-and-white house style: no colour, prints identically in mono.

\usepackage[T1]{fontenc}
\usepackage[utf8]{inputenc}
\usepackage{lmodern}
\usepackage[english]{babel}
\usepackage{microtype}
\usepackage[a4paper,margin=2.4cm]{geometry}
\usepackage{xcolor}
\usepackage{titlesec}
\usepackage{enumitem}
\usepackage{booktabs}
\usepackage{tabularx}
\usepackage{tcolorbox}
\tcbuselibrary{breakable,skins}
\usepackage{pifont}
\usepackage{xurl}
\usepackage[hidelinks]{hyperref}

\urlstyle{tt}

% --- Palette (plain black and white) ----------------------------------------
% Everything is black on white; the only greys are for a hairline rule and the
% callout title bar, so the documents print identically in mono.
\definecolor{sosAccent}{gray}{0}
\definecolor{sosDark}{gray}{0}
\definecolor{sosOpen}{gray}{0}
\definecolor{sosOpenBg}{gray}{1}
\definecolor{sosRule}{gray}{0.6}

% --- Headings ---------------------------------------------------------------
\titleformat{\section}{\normalfont\large\bfseries}{\thesection}{0.6em}{}
\titleformat{\subsection}{\normalfont\normalsize\bfseries\itshape}{\thesubsection}{0.6em}{}
\titlespacing*{\section}{0pt}{1.1\baselineskip}{0.35\baselineskip}
\titlespacing*{\subsection}{0pt}{0.7\baselineskip}{0.2\baselineskip}

\setlist[itemize]{leftmargin=1.2em,itemsep=0.15em,topsep=0.25em,parsep=0pt}
\setlist[enumerate]{leftmargin=1.4em,itemsep=0.15em,topsep=0.25em,parsep=0pt}

% --- Semantic helpers -------------------------------------------------------
% Inline reference to a file or path inside the project repository.
\newcommand{\repofile}[1]{\path{#1}}

% One "the assignment asks for ..." lead-in line.
\newcommand{\asks}[1]{\noindent\textbf{Assignment asks for:} #1\par\vspace{0.3em}}

% Open question the team still has to answer (was <mark>highlighted</mark> in the
% markdown draft).
\newenvironment{openbox}[1][Open --- needs the team]{%
  \begin{tcolorbox}[
    breakable,
    colback=white,
    colframe=black,
    colbacktitle=black!8,
    coltitle=black,
    boxrule=0.4pt,
    titlerule=0pt,
    left=0.6em,right=0.6em,top=0.5em,bottom=0.5em,
    title={#1},
    fonttitle=\bfseries\small,
    sharp corners,
  ]\small
}{%
  \end{tcolorbox}
}

\NewDocumentEnvironment{infobox}
  { O{Open --- needs the team} O{} }
{%
  \begin{tcolorbox}[
    breakable,
    colback=white,
    colframe=black,
    colbacktitle=black!8,
    coltitle=black,
    boxrule=0.4pt,
    titlerule=0pt,
    left=0.6em,
    right=0.6em,
    top=0.5em,
    bottom=0.5em,
    title={#1},
    fonttitle=\bfseries\small,
    sharp corners,
    #2
  ]\small
}{%
  \end{tcolorbox}
}


% Status markers for the submission checklist.
\newcommand{\stdone}{\ding{51}}
\newcommand{\stpartial}{\ding{108}}
\newcommand{\stgap}{\ding{55}}

% Title block.
\newcommand{\sostitle}[3]{%
  \begin{center}
    {\LARGE\bfseries #1}\\[0.35em]
    {\large\itshape #2}\\[0.5em]
    {\small #3}
  \end{center}
  \vspace{0.4em}
  {\color{sosRule}\hrule height 0.4pt}
  \vspace{0.9em}
}
.
% Plain black-and-white house style: no colour, prints identically in mono.

\usepackage[T1]{fontenc}
\usepackage[utf8]{inputenc}
\usepackage{lmodern}
\usepackage[english]{babel}
\usepackage{microtype}
\usepackage[a4paper,margin=2.4cm]{geometry}
\usepackage{xcolor}
\usepackage{titlesec}
\usepackage{enumitem}
\usepackage{booktabs}
\usepackage{tabularx}
\usepackage{tcolorbox}
\tcbuselibrary{breakable,skins}
\usepackage{pifont}
\usepackage{xurl}
\usepackage[hidelinks]{hyperref}

\urlstyle{tt}

% --- Palette (plain black and white) ----------------------------------------
% Everything is black on white; the only greys are for a hairline rule and the
% callout title bar, so the documents print identically in mono.
\definecolor{sosAccent}{gray}{0}
\definecolor{sosDark}{gray}{0}
\definecolor{sosOpen}{gray}{0}
\definecolor{sosOpenBg}{gray}{1}
\definecolor{sosRule}{gray}{0.6}

% --- Headings ---------------------------------------------------------------
\titleformat{\section}{\normalfont\large\bfseries}{\thesection}{0.6em}{}
\titleformat{\subsection}{\normalfont\normalsize\bfseries\itshape}{\thesubsection}{0.6em}{}
\titlespacing*{\section}{0pt}{1.1\baselineskip}{0.35\baselineskip}
\titlespacing*{\subsection}{0pt}{0.7\baselineskip}{0.2\baselineskip}

\setlist[itemize]{leftmargin=1.2em,itemsep=0.15em,topsep=0.25em,parsep=0pt}
\setlist[enumerate]{leftmargin=1.4em,itemsep=0.15em,topsep=0.25em,parsep=0pt}

% --- Semantic helpers -------------------------------------------------------
% Inline reference to a file or path inside the project repository.
\newcommand{\repofile}[1]{\path{#1}}

% One "the assignment asks for ..." lead-in line.
\newcommand{\asks}[1]{\noindent\textbf{Assignment asks for:} #1\par\vspace{0.3em}}

% Open question the team still has to answer (was <mark>highlighted</mark> in the
% markdown draft).
\newenvironment{openbox}[1][Open --- needs the team]{%
  \begin{tcolorbox}[
    breakable,
    colback=white,
    colframe=black,
    colbacktitle=black!8,
    coltitle=black,
    boxrule=0.4pt,
    titlerule=0pt,
    left=0.6em,right=0.6em,top=0.5em,bottom=0.5em,
    title={#1},
    fonttitle=\bfseries\small,
    sharp corners,
  ]\small
}{%
  \end{tcolorbox}
}

\NewDocumentEnvironment{infobox}
  { O{Open --- needs the team} O{} }
{%
  \begin{tcolorbox}[
    breakable,
    colback=white,
    colframe=black,
    colbacktitle=black!8,
    coltitle=black,
    boxrule=0.4pt,
    titlerule=0pt,
    left=0.6em,
    right=0.6em,
    top=0.5em,
    bottom=0.5em,
    title={#1},
    fonttitle=\bfseries\small,
    sharp corners,
    #2
  ]\small
}{%
  \end{tcolorbox}
}


% Status markers for the submission checklist.
\newcommand{\stdone}{\ding{51}}
\newcommand{\stpartial}{\ding{108}}
\newcommand{\stgap}{\ding{55}}

% Title block.
\newcommand{\sostitle}[3]{%
  \begin{center}
    {\LARGE\bfseries #1}\\[0.35em]
    {\large\itshape #2}\\[0.5em]
    {\small #3}
  \end{center}
  \vspace{0.4em}
  {\color{sosRule}\hrule height 0.4pt}
  \vspace{0.9em}
}


\hypersetup{pdftitle={Individual Reflection — Domi}}

\begin{document}

\sostitle{Individual Reflection}{Domi}{%
  StudyOS Thesis-Finder \quad\textbullet\quad individual submission \quad\textbullet\quad University of Tübingen}

\section*{1. Initial expectation}

Going into the course, I expected it to be more theory-heavy and less
project-driven, and I wasn't entirely sure what to expect content-wise. What I
did want was to deepen my skills in working with coding agents --- both learning
to work efficiently on my own with them and learning how to collaborate cleanly
with a team while using these tools. I was positively surprised by the module's
approach: in the first sessions we were forced to critically question what was
even worth ``solving,'' and how to approach the task at all, before writing a
line of code. The learning about the tools happened as a byproduct of actually
solving the assignment, not as a separate track. The one framework that stuck
with me is the double diamond: you have to question what is worth solving before
spending any thought on how to solve it or what to deliver --- because in the
age of coding agents, the ``how'' has almost become secondary, since so much
more is possible now that the real bottleneck moved upstream, to problem
selection. Initially I assumed we'd build a fully working ``assistant'' for
students that helps with everything, especially learning. That's technically
feasible, but it misses where the actual demand is.

\section*{2. Main learning}

My single most important learning is that, especially in the age of AI, the real
work is figuring out which problems are worth solving --- and that this is
exactly the kind of question that matters in product discovery and in startups:
where is there demand, where could demand emerge, and how do you even assess
that? Do other people actually share my framing of the problem, or am I
projecting my own view onto it? That's tied to actively seeking out other
people's opinions and talking to people who have relevant experience or a
valuable perspective --- a single viewpoint on a problem is almost never
complete. This really clicked for me when we pivoted from building an ``app'' to
building a ``skill.'' No user today wants to open a separate external tool just
to find a thesis topic. A lightweight skill that plugs the task directly into an
LLM the user already has open solves several of these problems at once. On top
of that, the simple exercise of writing ``we aim to solve X for Y by doing Z''
and testing it against impact, feasibility, and success already resolves a
surprising number of open questions and wrong turns before they cost any
implementation time.

\section*{3. Role and contribution in the group}

We were a team of three with fairly similar development roles overall, but each
of us leaned into a distinct strength. Max brought more expertise in building
out the app itself and in evaluating the skill's performance; Valentin focused
on outward communication and on user testing through his own network. My role
was largely a meta-role within that setup: repeatedly stepping back and asking
whether what we were currently doing actually made sense, and regularly calling
for a pause to make the current state explicit and to clarify the concrete next
steps for the group.

The clearest example of this dynamic was the no-DB pivot. The group was
initially divided on whether a database-backed or a database-free approach made
more sense --- I was the one pushing for no-DB. Rather than resolve it by
discussion alone, we pursued both in parallel for a short time. Once user
testing came back, it quickly became clear that the no-DB approach was not only
preferable from a user standpoint but also significantly less effort, especially
in terms of long-term maintenance. My most concrete individual contribution was
then following through on that direction: building the first no-DB prototype
myself, which ultimately became the final artifact we shipped.

\section*{4. Future learning}

Technically, I want to get better at working efficiently in larger AI-assisted
teams --- keeping communication clean, actually producing clean code, and not
letting errors creep in by holding the right balance between tests and
production code. I also want to keep sharpening my own ``vibe coding'' skill:
finding the right level of communication between me and the agent so that code
production runs as efficiently as possible, while still retaining enough detail
and understanding to stay in control of the product rather than losing oversight
of what's being built. Beyond that, working on this project noticeably grew my
interest in building my own products --- seeing what coding agents are actually
capable of, and getting a first real taste of how to design something with
genuine demand behind it rather than just something interesting to build, made
me want to pursue more of my own projects. To get there, I want to keep building
expertise independently so I can take on further personal projects in the near
future.

\end{document}
